\chapter{Полярные угловые блоки исправленного вихря}
\label{ch:v6-bosonic-defect-polar-angular-sturm-liouville-gate}

\section{Проблема}

Односторонняя декартова схема не сохраняет вращательную симметрию и потому
не может надёжно присвоить мягким уровням угловые числа. Требуется построить
гессиан непосредственно в полярных координатах, отделить переносную пару и
проверить внутренние каналы без решёточного смешения.

\section{Полярная вторая вариация}

В совместно вращающемся базисе возмущение с угловым числом $m$ записывается
через пять радиальных функций
\[
 (u,v,p,q,w)e^{im\theta},
\]
где $u,v$ --- продольная и фазовая компоненты заряженного поля, $p,q$ ---
радиальная и угловая компоненты связности, $w$ --- нейтральная амплитуда.
Использовано фоновое условие
\[
 \operatorname{div}\delta A-\frac{A}{G}(J\phi_0,\delta\phi)=0.
\]
Ему соответствует ковариантная переносная касательная $-D_j\phi$ с
компенсирующим калибровочным параметром $+A_j$.

В гессиан включены не только квадраты линейных вариаций $D\phi$, но и члены
второй вариации $-\delta A\,J\delta\phi$, сцепленные с ненулевой производной
фонового профиля. Их пропуск уничтожает точную переносную нулевую моду.

Условия регулярности в ядре наложены в круговом декартовом базисе:
\[
 u+iv,\ p+iq=O(r^{|m+1|}),\qquad
 u-iv,\ p-iq=O(r^{|m-1|}),\qquad
 w=O(r^{|m|}).
\]
Такая структура совпадает с частично-волновым разложением полного
калибровочно-скалярного оператора \cite{baacke-kevlishvili-partial-waves}
и с современной классификацией переносных и мультипольных мод
\cite{polar-alonso-vortex-spectrum}.

\section{Численный спектр}

Проверены угловые числа $-8\le m\le8$ на радиальных сетках $100$, $160$ и
$240$ узлов. Для $m=-1,0,1$ добавлена сетка $360$ узлов.

Нижний уровень каналов $m=\pm1$ равен
\[
 9.68349\cdot10^{-4},\quad
 3.95085\cdot10^{-4},\quad
 1.78475\cdot10^{-4},\quad
 7.99515\cdot10^{-5}.
\]
Порядок его сходимости к нулю равен $1.93674$. Следовательно, это именно две
вещественные переносные нулевые моды, представленные комплексными каналами
$m=\pm1$.

После удаления переносной пары минимальный внутренний уровень находится в
осесимметричном канале. На сетках $240$ и $360$ он равен
\[
 4.14560013,\qquad 4.14469858,
\]
а сдвиг составляет $9.02\cdot10^{-4}$. В остальных проверенных каналах
минимумы лежат выше $4.6989$ и возрастают к краям углового окна. Отрицательных
значений не найдено.

\begin{theorem}
В исправленной эффективной пятикомпонентной подсистеме вихря полярные блоки
$|m|\le8$ не содержат отрицательных мод. Единственные уровни, стремящиеся к
нулю, образуют переносную пару $m=\pm1$. Минимальный внутренний зазор в
проверенном окне сходится к значению около $4.1447$.
\end{theorem}

\section{Нулевая гипотеза и стоп-критерий}

Нулевая гипотеза следующего шага: центробежная часть делает все блоки
$|m|>8$ коэрцитивными, поэтому глобальный минимум внутреннего оператора уже
достигнут при $m=0$. Для закрытия требуется аналитическая нижняя оценка
высокоуглового хвоста. Простого наблюдения роста конечного набора уровней
недостаточно.

Стоп-критерий: если оценка хвоста допускает отрицательный блок или требует
нового свободного коэффициента, полный эффективный угловой сертификат не
проходит. Даже успешная оценка хвоста не заменяет гессиан всех компонент
$Q+T+B$.

\section{Статус}

Переносная пара строго разрешена, а проверенное окно $|m|\le8$ устойчиво.
Непрерывный хвост всех угловых чисел, полный спин-тензорный оператор,
замыкание нити и рождение наблюдаемой материи остаются открытыми.

\begin{thebibliography}{9}
\bibitem{baacke-kevlishvili-partial-waves}
J.~Baacke, N.~Kevlishvili,
\emph{One-loop corrections to the string tension of the vortex in the
Abelian Higgs model}, arXiv:0806.4349.

\bibitem{polar-alonso-vortex-spectrum}
A.~Alonso-Izquierdo, D.~Migu\'elez-Caballero, J.~Queiruga,
\emph{Spectral structure of fluctuations around $n$-vortices in the
Abelian-Higgs model}, arXiv:2505.05039.
\end{thebibliography}